% 微分同胚（综述）
% license CCBYSA3
% type Wiki

本文根据 CC-BY-SA 协议转载翻译自维基百科\href{https://en.wikipedia.org/wiki/Diffeomorphism}{相关文章}

在数学中，微分同胚是可微流形之间的一种同构。它是一个可逆函数，能够将一个可微流形映射到另一个可微流形，并且该函数及其逆函数都是连续可微的。
\begin{figure}[ht]
\centering
\includegraphics[width=6cm]{./figures/61cddbbc19414b20.png}
\caption{一个从正方形映射到其自身的微分同胚下，正方形上矩形网格的映射图像。} \label{fig_WeiFTP_1}
\end{figure}
\subsection{定义}
给定两个可微流形 $M$ 和 $N$，一个连续可微映射$f: M \to N$如果它是双射，并且它的逆映射$f^{-1}: N \to M$也是可微的，那么 $f$ 就称为一个微分同胚。如果映射$f$及其逆映射都是$r$ 次连续可微的，则称 $f$是一个$C^r$-微分同胚。

两个流形$M$和$N$如果存在一个微分同胚$f: M \to N$则称它们是微分同胚的，通常记作$M \simeq N$。如果两个流形都是$C^r$可微的，并且存在一个$r$次连续可微的双射，其逆映射也是$r$次连续可微的，那么称它们是$C^r$-微分同胚的。$C^1$-微分同胚通常直接称为微分同胚。$C^0$-微分同胚则等价于一个同胚。
\subsection{流形子集的微分同胚}
设$M$和$N$是两个流形，$X \subset M$ 和 $Y \subset N$ 是它们的子集。一个函数$f: X \to Y$称为光滑的，如果对每一个点 $p \in X$，都存在一个包含 $p$ 的邻域$U \subset M$以及一个光滑函数$g: U \to N$使得在交集部分有$g|_{U \cap X} = f|_{U \cap X}$。换句话说，函数 $g$ 是函数 $f$ 的一个光滑延拓。若函数 $f$ 是双射、光滑的，并且它的逆映射也是光滑的，那么 $f$ 就称为一个微分同胚。
\subsection{局部描述}
在一定的温和条件下，可以通过局部方法来检验一个可微映射是否是微分同胚。这就是 Hadamard–Caccioppoli 定理[1]：

若$U$和$V$是$\mathbb{R}^n$中的两个连通开子集，且$V$是单连通的，那么一个可微映射$f: U \to V$如果满足以下条件，则$f$是一个微分同胚：$f$是适当映射；对于每个$x \in U$，其微分$Df_x: \mathbb{R}^n \to \mathbb{R}^n$是双射（因此是一个线性同构）。

一些说明：

对于函数$f$来说，集合$V$的单连通性是保证$f$在全局上可逆的关键条件（仅依赖于导数在每一点都是双射这一条件）。

例如，考虑复平方函数的“实数化”版本：
  $$
  \begin{cases}
  f:\mathbb{R}^2 \setminus \{(0,0)\} \;\to\; \mathbb{R}^2 \setminus \{(0,0)\}\\
  (x,y) \;\mapsto\; (x^2 - y^2,\, 2xy).
  \end{cases}~
  $$
在这种情况下，函数 $f$ 是满射，并且满足：
  $$
  \det Df_x = 4(x^2 + y^2) \neq 0.~
  $$
也就是说，虽然在每一点上微分$Df_x$都是双射，但函数$f$并不是可逆的，因为它不是单射，例如：$f(1,0) = (1,0) = f(-1,0)$。

由于在可微函数中，一个点处的微分
  $$
  Df_x: T_xU \;\to\; T_{f(x)}V~
  $$
是一个线性映射，因此当且仅当$Df_x$是双射时，它才有定义良好的逆映射$Df_x$的矩阵形式是一个$n \times n$的矩阵，其第$i$行、第$j$列的元素是：$\partial f_i/\partial x_j$。这个矩阵被称为雅可比矩阵，它常用于显式计算。

微分同胚必须发生在相同维度的流形之间。设映射$f$从维度 $n$ 的流形映射到维度 $k$ 的流形：若$n < k$，则$Df_x$永远不可能是满射；若$n > k$，则$Df_x$永远不可能是单射。因此，在这两种情况下，$Df_x$都无法成为双射。

如果$Df_x$在某点$x$上是双射，则称$f$在$x$附近是一个局部微分同胚，因为由连续性可知，$x$附近所有足够接近的点$y$上的$Df_y$也都是双射。

若从维度$n$到维度$k$的光滑映射$f$满足$Df$（或局部的 $Df_x$）是满射，则称 $f$ 是一个沉降（局部称作“局部沉降”）。若$Df$（或局部的$Df_x$）是单射，则称$f$是一个浸入（局部称作“局部浸入”）。

一个可微的双射函数不一定是微分同胚。例如：$f(x) = x^3$并不是 $\mathbb{R}$ 到 $\mathbb{R}$ 的微分同胚，因为它的导数在 $0$ 处为零，从而它的逆函数在 $0$ 处不可微。这是一个同胚但不是微分同胚的例子。

当$f$是可微流形之间的映射时，“微分同胚”比“同胚”是更强的条件：对于微分同胚，$f$及其逆映射都必须是可微的；对于同胚，$f$及其逆映射只需要连续即可。因此，每个微分同胚都是同胚，但并非每个同胚都是微分同胚。

一个映射$f: M \to N$是微分同胚，当且仅当它在局部坐标图下满足微分同胚的定义。更具体地：为$M$选取一组相容的坐标图，为$N$也选取一组相容的坐标图。令$\phi$ 是$M$上的坐标图，$\psi$是$N$上的坐标图，对应的像分别为$U$和$V$。若对所有满足$f(\phi^{-1}(U)) \subseteq \psi^{-1}(V)$的情况，组合映射$\psi f \phi^{-1}: U \to V$是一个微分同胚，则$f$是$M$到$N$的微分同胚。
\subsection{示例}
由于任何流形都可以在局部参数化，我们可以考虑一些从$\mathbb{R}^2 \to \mathbb{R}^2$的显式映射。
\begin{itemize}
\item 设：
$$
f(x, y) = \bigl(x^2 + y^3,\; x^2 - y^3 \bigr)~
$$
我们计算它的雅可比矩阵：
$$
J_f = 
\begin{pmatrix}
2x & 3y^2 \\
2x & -3y^2
\end{pmatrix}~
$$
该矩阵的行列式为零，当且仅当：$xy = 0$。这意味着，只有在远离$x$-轴和$y$-轴的区域，$f$才有可能是局部微分同胚。然而，$f$并不是双射，因为：$f(x, y) = f(-x, y)$，因此它不能是一个微分同胚。
\item 设：
$$
g(x, y) =
\bigl(
a_{0} + a_{1,0}x + a_{0,1}y + \cdots ,\;
b_{0} + b_{1,0}x + b_{0,1}y + \cdots
\bigr)~
$$
其中$a_{i,j}$和$b_{i,j}$是任意实数，省略的项都是关于$x$和$y$至少二次的高阶项。在$(0,0)$处的雅可比矩阵为：
$$
J_{g}(0,0) =
\begin{pmatrix}
a_{1,0} & a_{0,1} \\
b_{1,0} & b_{0,1}
\end{pmatrix}.~
$$
可以看出，当且仅当：
$$
a_{1,0} b_{0,1} - a_{0,1} b_{1,0} \neq 0~
$$
时，$g$在$(0,0)$处是一个局部微分同胚。换句话说，$g$的分量函数的线性项作为多项式必须线性无关。
\item 设：
$$
h(x, y) = \bigl(\sin(x^2 + y^2),\; \cos(x^2 + y^2)\bigr)~
$$
计算其雅可比矩阵：
$$
J_h =
\begin{pmatrix}
2x\cos(x^2 + y^2) & 2y\cos(x^2 + y^2)\\
-2x\sin(x^2 + y^2) & -2y\sin(x^2 + y^2)
\end{pmatrix}.~
$$
可以看到，这个雅可比矩阵的行列式在所有点都为零！事实上，$h$的像仅为单位圆：$\sin^2(x^2 + y^2) + \cos^2(x^2 + y^2) = 1$.
\end{itemize}
\subsubsection{表面形变}
在力学中，由应力引起的变换称为形变，这种变换可以用一个微分同胚来描述。设$f: U \to V$是两个表面$U$和$V$之间的一个微分同胚，那么其雅可比矩阵$Df$是一个可逆矩阵。事实上，要求对$U$中的每一点 $p$，其某个邻域内的雅可比矩阵$Df$都保持非奇异。假设在某个表面的坐标图下：$f(x, y) = (u, v)$.

全微分$u$的全微分为：
$$
du = \frac{\partial u}{\partial x} dx + \frac{\partial u}{\partial y} dy~
$$
同理，$v$的全微分也有类似表达式。

于是：$(du, dv) = (dx, dy) Df$,这表示一个线性变换，它以原点为固定点，可以看作是某种类型复数的作用。当$(dx, dy)$也被解释为同一类型的复数时，这个线性作用就对应于该复数平面中的复数乘法。由于复数乘法保留某种类型的角度（欧几里得角、双曲角或斜率角），且$Df$在表面上始终可逆，因此该类型的复数在整个表面上是一致的。由此，一个表面形变或两个表面之间的微分同胚具有一种保角性，即保留某种意义下的角度。
\subsection{微分同胚群}
设$M$是一个可微流形，并且它是第二可数且豪斯多夫的。$M$的微分同胚群是指所有从 $M$到$M$的$C^r$微分同胚所组成的群，记作：$\text{Diff}^r(M)$当$r$已知或不加特别说明时，通常简写为：$\text{Diff}(M)$。这是一个“庞大”的群，因为只要 $M$不是零维流形，这个群就不是局部紧致的。
\subsubsection{拓扑结构}
微分同胚群有两种自然的拓扑：弱拓扑和强拓扑（Hirsch，1997）。当流形是紧致的，这两种拓扑是一致的。弱拓扑总是可度量化的。当流形是非紧致的，强拓扑可以刻画函数在“无穷远处”的行为，但它不可度量化。不过，它依然是一个Baire空间。

在流形 $M$ 上固定一个黎曼度量，弱拓扑由如下度量族诱导：
$$
d_K(f,g) = \sup_{x \in K} d(f(x), g(x)) + 
\sum_{1 \le p \le r} \sup_{x \in K} \|D^p f(x) - D^p g(x)\|,~
$$
其中 $K$ 是 $M$ 的紧子集，$\|D^p f(x) - D^p g(x)\|$ 表示 $p$ 阶导数的差值的范数。

由于 $M$ 是 $\sigma$-紧致的（$\sigma$-compact），存在一列紧集 $\{K_n\}$，其并集为 $M$。于是可以定义全局的距离：
$$
d(f,g) = \sum_n 2^{-n} \cdot 
\frac{d_{K_n}(f,g)}{1 + d_{K_n}(f,g)}~
$$
赋予弱拓扑后，微分同胚群在局部上与$C^r$向量场的空间同胚（Leslie，1967）。
在紧子集上，通过固定一个黎曼度量并使用该度量的指数映射，可以将微分同胚群局部识别为向量场的空间。如果$r$是有限的且流形是紧致的，那么向量场空间是一个巴拿赫空间。此时，不同图表之间的过渡映射是光滑的，使得微分同胚群成为一个具有平滑右平移的巴拿赫流形；不过左平移和求逆操作仅是连续的。如果$r = \infty$，那么向量场空间是一个弗雷歇空间，且过渡映射是光滑的，使微分同胚群成为一个弗雷歇流形，甚至是一个正规弗雷歇李群。如果流形是$\sigma$-紧致但非紧致，则在任意这两种拓扑下，整个微分同胚群都不是局部可收缩的。此时，需要通过控制群在“无穷远处”的偏离程度来限制群的范围，才能得到一个具有流形结构的微分同胚群（参见 Michor & Mumford 2013）。
\subsubsection{李代数}
流形$M$的微分同胚群的李代数由$M$上的所有向量场组成，并配备向量场的李括号作为运算。从形式上看，这可以通过在空间中每个点的坐标$x$做一个微小变化来理解：
$$
x^\mu \;\mapsto\; x^\mu + \varepsilon h^\mu(x),~
$$
其中$\varepsilon$是一个无穷小参数，$h^\mu(x)$是一个光滑函数。对应的无穷小生成元就是向量场：
$$
L_h = h^\mu(x) \frac{\partial}{\partial x^\mu}~
$$
\subsubsection{示例}
\begin{itemize}
\item 当$M = G$是一个李群时，$G$可以通过左平移自然地嵌入到其自身的微分同胚群中。设$\text{Diff}(G)$表示 $G$ 的微分同胚群，则存在如下分解：$\text{Diff}(G) \;\simeq\; G \times \text{Diff}(G, e)$，其中$\text{Diff}(G, e)$是在$G$的单位元$e$处固定不动的微分同胚子群。
\item 欧几里得空间$\mathbb{R}^n$的微分同胚群有两个连通分支：保定向的微分同胚；反转定向的微分同胚。实际上，一般线性群 $\text{GL}(n, \mathbb{R})$ 是如下子群$
\text{Diff}(\mathbb{R}^n, 0)$的形变收缩，这里 $\text{Diff}(\mathbb{R}^n, 0)$ 表示所有在原点处固定不动的微分同胚。该收缩由以下映射定义：$f(x) \;\mapsto\; f(tx)/t, \quad t \in (0, 1]$。因此，一般线性群也是整个微分同胚群的一个形变收缩。
\item 对于一个有限点集，其微分同胚群就是对称群。类似地，如果 $M$ 是任意流形，则存在如下群扩张：$0 \;\longrightarrow\; \text{Diff}_0(M)\;\longrightarrow\; \text{Diff}(M)\;\longrightarrow\; \Sigma(\pi_0(M))$其中：$\text{Diff}_0(M)$ 是保持流形所有连通分支的微分同胚子群；
$\Sigma(\pi_0(M))$ 是连通分支集合 $\pi_0(M)$ 的排列群。此外，映射$\text{Diff}(M) \;\to\; \Sigma(\pi_0(M))$的像是那些保持各连通分支微分同胚类的双射。
\end{itemize}
\subsubsection{传递性}
对于一个连通流形$M$，其微分同胚群在$M$上是传递作用的。更一般地，微分同胚群在构型空间$C_kM$上也是传递作用的。如果$M$至少是二维的，则微分同胚群在有序构型空间$F_kM$上也是传递作用的，并且在$M$本身上的作用是多重传递的（Banyaga 1997，第29页）。
\subsubsection{微分同胚的扩展}
1926年，Tibor Radó 提出一个问题：单位圆上的任意同胚或微分同胚的调和扩展是否会在开单位圆盘上产生一个微分同胚。随后不久，Hellmuth Kneser给出了一个优雅的证明。1945年，Gustave Choquet显然并不知道这一结果，又独立地给出了一个完全不同的证明。

圆的（保持定向的）微分同胚群是路径连通的。这可以通过注意到任意这样的微分同胚都可以提升为实数上的一个微分同胚$f$，且满足：$[f(x+1) = f(x) + 1]$,而该空间是凸的，因此是路径连通的。通过一条最终收敛于恒等映射的光滑路径，还可以用另一种更为基础的方法将圆上的微分同胚扩展到开单位圆盘上（这是一种 Alexander trick 的特例）。此外，圆的微分同胚群具有与正交群 $O(2)$ 相同的同伦型。

在1950和1960年代，高维球面$S^{n-1}$上微分同胚的类似扩展问题得到了广泛研究，René Thom、John Milnor和Stephen Smale都做出了重要贡献。这类扩展的障碍由一个有限的阿贝尔群 $\Gamma_n$（称为“扭曲球面的群”）给出，它被定义为微分同胚群的阿贝尔分量群与可以扩展为球体 $B^n$ 微分同胚的类所构成的子群的商。
\subsubsection{连通性}
对于流形而言，微分同胚群通常不是连通的。它的分量群被称为映射类群。在二维情形（即曲面）中，映射类群是一个有限呈示群，由Dehn扭转生成；这一结果由Max Dehn、W. B. R. Lickorish和Allen Hatcher证明。Max Dehn和Jakob Nielsen还证明了该群可以被识别为曲面基本群的外自同构群。

William Thurston 对这一分析进行了改进，将映射类群的元素分为三类：与某个周期性微分同胚等价的元素；与保持某条简单闭曲线不变的微分同胚等价的元素；以及与伪 Anosov 微分同胚等价的元素。在环面$S^1 \times S^1 = \mathbb{R}^2 / \mathbb{Z}^2$的情形下，映射类群就是模群$\text{SL}(2, \mathbb{Z})$,此时的分类在椭圆型、抛物型和双曲型矩阵的意义下变成了经典理论。

Thurston通过观察映射类群自然地作用于Teichmüller空间的紧化空间完成了这一分类；由于该扩展空间与一个闭球同胚，因此Brouwer不动点定理可以适用。Smale猜想，如果$M$是一个定向的光滑闭流形，则保持定向的微分同胚群的恒等分支是单纯群。这一猜想首先由Michel Herman在圆环乘积的情形下得到了证明；Thurston则在更一般的情形下给出了完整的证明。
\subsubsection{同伦类型}
\begin{itemize}
\item $S^2$ 的微分同胚群具有与其子群 $\mathrm{O}(3)$ 相同的同伦类型。这一点由 Steve Smale 证明。[2]
\item 环面的微分同胚群具有与其线性自同构群相同的同伦类型：$S^1 \times S^1 \times \mathrm{GL}(2, \mathbb{Z})$.
\item 属数$g > 1$的可定向曲面的微分同胚群具有与其映射类群相同的同伦类型（即各个分支都是可收缩的）。
\item 三维流形的微分同胚群的同伦类型已经通过 Ivanov、Hatcher、Gabai和 Rubinstein的工作得到了相当深入的理解，尽管仍有一些未解决的问题，主要集中在基本群有限的三维流形上。
\item 对于$n > 3$的$n$维流形，其微分同胚群的同伦类型目前了解得并不充分。例如，关于$\mathrm{Diff}(S^4)$是否具有超过两个连通分支仍然是一个悬而未决的问题。然而，通过 Milnor、Kahn 和 Antonelli 的研究，已经知道当$n > 6$时，$\mathrm{Diff}(S^n)$并不具有有限 CW 复形的同伦类型。
\end{itemize}
**Homeomorphism and diffeomorphism（同胚与微分同胚）**

由于每个微分同胚都是一个同胚，如果一对流形彼此微分同胚，那么它们特别地也是同胚的。然而，反过来一般并不成立。

虽然很容易找到一些是同胚但不是微分同胚的映射，但要找到一对彼此同胚却不是微分同胚的流形则更困难。在 1、2 和 3 维情形中，任意一对同胚的光滑流形都是微分同胚的。而在 4 维或更高维度中，存在同胚但不微分同胚的流形对。第一个这样的例子由 John Milnor 在 7 维情形下构造出来。他构造了一个光滑的 7 维流形（现在称为 Milnor 球），它与标准的 7 球是同胚的，但不是微分同胚的。实际上，与 7 球同胚的流形存在 28 种有向微分同胚类，每一种都是一个以 3 球为纤维、4 球为底空间的纤维丛的总空间。

在 4 维流形中，还出现了更奇特的现象。20 世纪 80 年代初，Simon Donaldson 和 Michael Freedman 的成果结合起来，发现了所谓的**异类 $\mathbb{R}^4$**：存在不可数多个两两不微分同胚的 $\mathbb{R}^4$ 开子集，每一个都与标准的 $\mathbb{R}^4$ 同胚；此外，还存在不可数多个两两不微分同胚的、与 $\mathbb{R}^4$ 同胚但无法光滑嵌入 $\mathbb{R}^4$ 的可微流形。

**See also**

* Anosov 微分同胚，例如 Arnold 的猫映射
* Diffeo anomaly（微分同胚异常），也称为引力异常，是量子力学中的一种异常
* Diffeology，集合上的光滑参数化，形成一个微分同胚空间
* Diffeomorphometry，计算解剖学中关于形状和形式的度量研究
* Étale morphism
* Large diffeomorphism（大微分同胚）
* Local diffeomorphism（局部微分同胚）
* Superdiffeomorphism（超微分同胚）

**Notes**

* Steven G. Krantz; Harold R. Parks (2013). *The implicit function theorem: history, theory, and applications*. Springer. p. Theorem 6.2.4. ISBN 978-1-4614-5980-4.
* Smale (1959). "Diffeomorphisms of the 2-sphere". *Proc. Amer. Math. Soc.* 10 (4): 621–626. doi:10.1090/s0002-9939-1959-0112149-8.
